% cap4_implementare.tex
\chapter{Implementarea prototipului}
\label{ch:implementare}

Capitolul de față descrie concret modul în care componentele prezentate în
\chapref{ch:metodologie} au fost transpuse în cod funcțional și asamblate
într-un prototip operațional.
Sunt tratate, în ordine, firmware-ul nodului interior bazat pe
\gls{esp32}, firmware-ul camerei exterioare \gls{esp32}-CAM, configurarea
gateway-ului Raspberry Pi, structura \gls{api}-ului NestJS, motorul de
detecție a anomaliilor cu \gls{if} și, în final, aplicația mobilă.

% ──────────────────────────────────────────────────────────────────
\section{Firmware-ul nodului interior (ESP32 WROOM-32)}
\label{sec:firmware-nod}

\subsection{Pinout și configurare hardware}
\label{sec:pinout}

Nodul interior este construit în jurul modulului ESP32 WROOM-32.
Toți pinii \gls{adc} folosiți pentru senzori (MQ-2, \gls{ldr}\,\#1 și
\gls{ldr}\,\#2) aparțin grupului ADC1 (GPIO34, GPIO35, GPIO36), ceea ce
permite citiri analogice corecte chiar și cu \gls{wifi} activ — spre
deosebire de pinii ADC2 care devin inutilizabili în prezența
radio-ului~\cite{espressif2023esp32}.
Modulul BH1750 (citire de lux prin I\textsuperscript{2}C), luat în
considerare inițial, a fost eliminat în favoarea celor două \gls{ldr}
analogice, simplificând schema de cablare.

\tabref{tab:pinout} centralizează cei unsprezece pini activi ai
nodului, împărțiți în senzori și actuatori.

\begin{table}[H]
  \centering
  \caption{Pinout complet al nodului interior (ESP32 WROOM-32).}
  \label{tab:pinout}
  \begin{tabular}{@{} l l p{6.2cm} @{}}
    \toprule
    \textbf{Componentă} & \textbf{Pin GPIO} & \textbf{Observații} \\
    \midrule
    \gls{dht}11 \#1 (living)   & GPIO4  & Digital, pull-up extern 10\,k\textOmega{} la 3{,}3\,V \\
    \gls{dht}11 \#2 (dormitor) & GPIO23 & Digital, pull-up extern 10\,k\textOmega{} la 3{,}3\,V \\
    MQ-2 (gaz, A0)             & GPIO34 & ADC1, intrare; divizor de tensiune 0–3{,}3\,V \\
    \gls{ldr} \#1 (lumină)     & GPIO35 & ADC1, intrare; divizor cu 10\,k\textOmega{} \\
    \gls{ldr} \#2 (lumină)     & GPIO36 & ADC1 (VP), intrare; divizor cu 10\,k\textOmega{} \\
    Buzzer activ               & GPIO5  & Ieșire digitală, oscilator intern \\
    Servo SG90 (draperii)       & GPIO18 & \gls{pwm} 50\,Hz; alimentat din 5\,V extern \\
    \gls{led} \#0 (living)     & GPIO25 & Ieșire digitală + rezistor serie 220\,\textOmega{} \\
    \gls{led} \#1 (dormitor)   & GPIO26 & Ieșire digitală + rezistor serie 220\,\textOmega{} \\
    \gls{led} \#2 (hol)        & GPIO27 & Ieșire digitală + rezistor serie 220\,\textOmega{} \\
    Ventilator 5\,V            & GPIO33 & Comandat prin tranzistor NPN (NU direct din GPIO) \\
    \bottomrule
  \end{tabular}
\end{table}

\figref{fig:pinout} ilustrează schema de conexiuni a nodului, cu marcarea
alimentărilor separate: 3{,}3\,V pentru senzori (\gls{dht}11, \gls{ldr})
și 5\,V extern pentru servo și ventilator.

\begin{figure}[H]
  \centering
  \includegraphics[width=\singlefigure]{\pics/fig-pinout}
  \caption{Schema de conexiuni a nodului interior ESP32 WROOM-32
           (pinout conform \tabref{tab:pinout}).}
  \label{fig:pinout}
\end{figure}

\subsection{Citirea senzorilor}
\label{sec:citire-senzori}

La inițializare, rezoluția \gls{adc} este setată explicit la 10 biți
(\texttt{analogReadResolution(10)}), oferind valori în intervalul 0–1023.
Ambii senzori \gls{dht}11 sunt citiți cu biblioteca \emph{Adafruit DHT},
iar erorile temporare (valori \texttt{NaN}) sunt tratate tolerant:
dacă cel puțin un senzor raportează date valide, publicarea continuă cu
câmpurile disponibile; numai dacă ambii eșuează simultan este omisă
întreaga publicare~\cite{aosong2019dht11}.

Senzorul de gaz MQ-2 returnează o valoare analogică brută (0–1023).
Când valoarea depășește pragul \texttt{GAS\_THRESHOLD = 350}, firmware-ul
detectează \emph{frontul crescător} (tranziția \texttt{false → true}),
publică o alertă pe topicul \texttt{smarthome/alerts} și activează
buzzerul local cu trei bipuri — prevenind astfel spamul la broker în caz
de depășire prelungită~\cite{hanwei2020mq2}.
\listingref{lst:publish-sensors} prezintă nucleul funcției de colectare
și publicare.

\begin{lstlisting}[language=C++,
  caption={Fragment din \texttt{publishSensors()} — citire DHT11, MQ-2,
           LDR și construire payload JSON (esp32\_node.ino).},
  label=lst:publish-sensors]
// Citire senzori ADC — 10 biti (0-1023)
int gas    = analogRead(MQ2_PIN);   // GPIO34
int light1 = analogRead(LDR1_PIN);  // GPIO35
int light2 = analogRead(LDR2_PIN);  // GPIO36

// Alerta gaz — doar pe frontul crescator
bool gasAlert = (gas > GAS_THRESHOLD);
if (gasAlert && !gasAlertActive) {
  publishGasAlert(gas);
  beep(3);         // semnal sonor local imediat
}
gasAlertActive = gasAlert;

// Payload JSON (contract cu backend si mobile)
StaticJsonDocument<384> doc;
doc["nodeId"]    = NODE_ID;          // "esp32_node_a"
doc["location"]  = LOCATION;        // "interior"
if (dht1ok) { doc["temperature"] = temp1; doc["humidity"] = humid1; }
if (dht2ok) { doc["temperature2"]= temp2; doc["humidity2"]= humid2; }
doc["gasLevel"]  = gas;
doc["gasAlert"]  = gasAlert;
doc["light1"]    = light1;
doc["light2"]    = light2;
doc["lightLux"]  = light1;          // alias backward-compat
\end{lstlisting}

\subsection{Actuatori}
\label{sec:actuatori}

Ventilatorul (5\,V) nu poate fi alimentat direct din \gls{gpio} —
curentul maxim al unui pin ESP32 este de 12\,mA — astfel că este comandat
printr-un tranzistor NPN pe GPIO33: nivel HIGH pe pin conduce tranzistorul
și închide circuitul motorului.
Servo-ul SG90, folosit pentru controlul draperiilor, primește semnal
\gls{pwm} la 50\,Hz prin biblioteca \emph{ESP32Servo}; funcția
\texttt{servoDraperii.write(angle)} acceptă unghiuri între 0° și 180°
(corespunzând draperiilor complet deschise, respectiv complet închise).

Ventilatorul poate fi acționat în două moduri, comutabile din aplicație:
\emph{manual} (comenzile \texttt{fan\_on}/\texttt{fan\_off}) și
\emph{automat} (comanda \texttt{fan\_auto}). În modul automat, firmware-ul
pornește ventilatorul când temperatura medie a celor doi senzori \gls{dht}11
depășește un \textbf{prag configurabil de la distanță} (comanda
\texttt{fan\_threshold}, implicit 28\,°C) și îl oprește sub
(prag~$-$~2\,°C), histerezisul reducând oscilațiile rapide în jurul pragului.
Atât modul, cât și pragul sunt salvate în memoria NVS și raportate înapoi în
mesajul de senzori (\texttt{fanOn}, \texttt{fanAuto}, \texttt{fanThreshold}),
astfel încât aplicația reflectă starea reală a dispozitivului.

\subsection{Comunicație MQTT, provisioning Wi-Fi și mDNS}
\label{sec:mqtt-provisioning}

La pornire, nodul citește credențialele Wi-Fi și adresa brokerului
\gls{mqtt} din memoria NVS (\emph{Non-Volatile Storage}) folosind
biblioteca \emph{Preferences}.
Dacă rețeaua salvată nu răspunde, dispozitivul revine automat la
hotspot-ul de configurare (\texttt{SmartHome-Setup}), de unde o comandă
\texttt{wifi\_update} — primită ulterior tot prin \gls{mqtt} — poate
instala noi credențiale și reporni nodul.

Adresa brokerului poate fi specificată ca nume mDNS (\texttt{smarthome.local});
în acest caz, funcția \texttt{resolveMqttHost()} interogează serverul mDNS
local (până la cinci tentative) și returnează IP-ul corespunzător,
eliminând necesitatea configurării unui IP static pe gateway~\cite{rfc6762mdns}.

Conexiunea \gls{mqtt} este stabilită cu autentificare (utilizator și parolă)
și cu un mesaj \emph{Last Will and Testament} reținut pe topicul
\texttt{smarthome\slb{} esp32\_node\_a\slb{}status}:
dacă brokerul detectectează o deconectare bruscă, publică automat
\texttt{\{"node\_id":\allowbreak"esp32\_node\_a",\allowbreak"status":\allowbreak"offline"\}}, informând
aplicația mobilă în timp real~\cite{oasis2019mqtt}.
La conectare cu succes, nodul se abonează la
\texttt{smarthome\slb{} esp32\_node\_a\slb{}commands} și publică imediat statusul
\texttt{online}.
Publicarea datelor de senzori are loc la fiecare 5\,s pe topicul
\texttt{smarthome\slb{} esp32\_node\_a\slb{}sensors}.

% ──────────────────────────────────────────────────────────────────
\section{Firmware-ul ESP32-CAM (nod exterior)}
\label{sec:firmware-cam}

Nodul exterior este realizat pe un modul AI-Thinker ESP32-CAM echipat cu
senzorul de imagine OV2640~\cite{espressif2023esp32cam}.
Singurul senzor de mișcare din sistem — un \gls{pir} HC-SR501 — este
conectat la GPIO13~\cite{pir2019hcsr501}.
La detecția mișcării, un \gls{led} exterior conectat la GPIO12 se aprinde
timp de 30\,s (gestionat prin variabila \texttt{lastMotion} și temporizare
cu \texttt{millis()}).

\subsection{Serverul HTTP și stream-ul MJPEG}
\label{sec:mjpeg}

Camera este inițializată prin \texttt{esp\_camera\_init()}, cu rezoluție
SVGA (800$\times$600) dacă este prezentă memoria PSRAM, sau VGA în lipsa
acesteia.
Firmware-ul pornește un server HTTP pe portul 80 cu două rute:
\begin{itemize}
  \item \texttt{GET /capture} — returnează un snapshot JPEG al cadrului
        curent (\texttt{Content-Type: image/jpeg});
  \item \texttt{GET /stream} — transmite un flux \emph{Motion JPEG}
        (\texttt{multipart/x-mixed-replace}), cu o pauză de 50\,ms între
        cadre ($\approx$15–20 fps), lăsând procesorului timp și pentru
        bucla \gls{mqtt}.
\end{itemize}

Camera devine accesibilă în rețeaua locală ca
\texttt{http://smarthome-cam.local} prin înregistrarea unui serviciu mDNS
(\texttt{MDNS.begin("smarthome-cam")}), similar nodului interior.
Tot periodic (la 5\,s), nodul publică un mesaj JSON pe topicul
\texttt{smarthome/esp32\_cam\_node/sensors} conținând câmpul \texttt{motion}
(0/1), iar conexiunea \gls{mqtt} folosește același mecanism
\emph{Last Will} ca nodul interior.

% ──────────────────────────────────────────────────────────────────
\section{Gateway Raspberry Pi}
\label{sec:gateway}

Gateway-ul rulează pe un Raspberry Pi 3B+ și găzduiește toate serviciile
de infrastructură ale sistemului.

\textbf{Mosquitto} este brokerul \gls{mqtt} (portul 1883), configurat cu
autentificare prin fișier de parole.
Acesta primește publicările nodurilor ESP32, le distribuie abonaților
(\gls{api}-ului NestJS și motorului ML) și gestionează mesajele
\emph{retained} pentru statusul nodurilor.

\textbf{PostgreSQL + TimescaleDB} stochează toate citirile de senzori în
hypertabela \texttt{sensor\_readings}, optimizată pentru interogări
time-series prin indexuri pe coloana \texttt{time}~\cite{timescaledb2017,postgresql2024}.
Celelalte tabele ale schemei (\texttt{alerts}, \texttt{nodes},
\texttt{automation\_rules}, \texttt{users}, \texttt{ml\_models}) sunt
tabele PostgreSQL standard.

Fiecare componentă software este înregistrată ca serviciu \textbf{systemd},
asigurând pornirea automată la boot și repornirea în caz de eroare:
\texttt{smarthome-api.service} (NestJS),
\texttt{smarthome-ml.service} (antrenare zilnică la 03:00) și
\texttt{mosquitto.service}.

% ──────────────────────────────────────────────────────────────────
\section{API NestJS}
\label{sec:api}

\gls{api}-ul backend este implementat cu framework-ul NestJS (Node.js,
TypeScript) și rulează pe portul 3000 al gateway-ului~\cite{nestjs2024}.
TypeORM asigură maparea obiect-relațională față de PostgreSQL.

\subsection{Structura modulelor}
\label{sec:module-nestjs}

\texttt{AppModule} importă nouă module funcționale:
\texttt{AuthModule}, \texttt{UsersModule}, \texttt{SensorsModule},
\texttt{AlertsModule}, \texttt{RulesModule}, \texttt{CommandsModule},
\texttt{MqttModule}, \texttt{SimulatorModule} și \texttt{SetupModule}.
Comunicarea inter-modul se realizează prin \texttt{EventEmitter2}
(biblioteca \emph{@nestjs/event-emitter}), evitând dependențe circulare.

\subsection{Subscriber MQTT}
\label{sec:subscriber-mqtt}

\texttt{MqttService} implementează \texttt{OnModuleInit} și se conectează
la Mosquitto la pornire (dezactivat automat dacă \texttt{SIMULATE=true}).
Se abonează la trei familii de topicuri cu wildcard:
\texttt{smarthome/+/sensors} (date senzori), \texttt{smarthome/alerts}
(alerte hardware) și \texttt{smarthome/+/status} (prezența nodurilor).
\listingref{lst:mqtt-handler} prezintă logica de procesare a mesajelor.

\begin{lstlisting}[language=TypeScript,
  caption={Handler de mesaje MQTT în \texttt{MqttService}
           (\texttt{mqtt.service.ts}) — ramificarea pe tip de topic.},
  label=lst:mqtt-handler]
private async handleMessage(topic: string, payload: string) {
  const data = JSON.parse(payload);

  if (topic.endsWith('/sensors')) {
    await this.sensorsService.saveReading({
      nodeId:       data.nodeId ?? data.node_id,
      temperature:  data.temperature,
      humidity:     data.humidity,
      temperature2: data.temperature2 ?? null,
      humidity2:    data.humidity2    ?? null,
      gasLevel:     data.gasLevel,
      gasAlert:     data.gasAlert ?? false,
      light1:       data.light1   ?? null,
      light2:       data.light2   ?? null,
      lightLux:     data.lightLux ?? data.light1,
      motion:       data.motion   ?? false,
    });
  } else if (topic === 'smarthome/alerts') {
    await this.alertsService.create({
      nodeId:    data.node_id,
      alertType: data.alert_type,   // ex. "GAS_DETECTED"
      severity:  data.alert_type === 'GAS_DETECTED' ? 'critical' : 'medium',
      location:  data.location,
      details:   data.details,
    });
  } else if (topic.endsWith('/status')) {
    await this.sensorsService.updateNodeStatus(data);
  }
}
\end{lstlisting}

Publicarea comenzilor spre noduri se face prin evenimentul
\texttt{command.sent} emis de \texttt{CommandsService} și capturat de
\texttt{MqttService}, care îl publică pe topicul
\texttt{smarthome/\{nodeId\}/commands}.

\subsection{WebSocket (Socket.IO)}
\label{sec:websocket}

\texttt{SensorsGateway} declară un namespace Socket.IO \texttt{/live} și
reacționează la trei evenimente interne:
\texttt{sensor.reading} → emite \texttt{sensor\_update} (datele noi de senzori),
\texttt{alert.new} → emite \texttt{alert},
\texttt{node.status} → emite \texttt{node\_status}.
La conectarea unui client, tokenul \gls{jwt} este verificat opțional;
dacă tokenul e prezent dar invalid, clientul este deconectat imediat.

\subsection{Endpointuri REST și autentificare JWT}
\label{sec:rest}

Toate rutele protejate sunt acoperite de garda \texttt{JwtAuthGuard}.
Tabelul următor rezumă principalele endpointuri~\cite{jones2015jwt}:

\begin{table}[H]
  \centering
  \caption{Principalele endpointuri \gls{rest} ale \gls{api}-ului NestJS.}
  \label{tab:rest-endpoints}
  \begin{tabular}{@{} l l p{5.5cm} @{}}
    \toprule
    \textbf{Metodă} & \textbf{Cale} & \textbf{Descriere} \\
    \midrule
    POST & \texttt{/api/auth/login}             & Autentificare; returnează token JWT \\
    GET  & \texttt{/api/sensors/latest}         & Ultima citire per nod \\
    GET  & \texttt{/api/sensors/nodes}          & Lista nodurilor cu status online \\
    GET  & \texttt{/api/sensors/history}        & Istoric cu filtre \texttt{nodeId}, \texttt{from}, \texttt{to}, \texttt{limit} \\
    GET  & \texttt{/api/alerts}                 & Toate alertele \\
    POST & \texttt{/api/alerts/\{id\}/acknowledge} & Marchează alertă ca recunoscută \\
    POST & \texttt{/api/commands/\{nodeId\}}    & Trimite comandă actuator \\
    GET  & \texttt{/api/rules}                  & Reguli de automatizare \\
    GET  & \texttt{/api/setup/status}           & Status provisionare Wi-Fi \\
    \bottomrule
  \end{tabular}
\end{table}

Autentificarea utilizează JWT cu semnătură HS256.
Serviciul \texttt{AuthService} verifică parola prin \texttt{bcrypt.compare()}
față de hash-ul stocat în baza de date~\cite{provos2003bcrypt},
iar \texttt{JwtService.sign()} emite tokenul cu payload
\texttt{\{sub, username\}}.

\subsection{Modulul simulator}
\label{sec:simulator}

Pentru dezvoltare fără hardware fizic, variabila de mediu
\texttt{SIMULATE=true} activează \texttt{SimulatorService}.
Acesta generează la fiecare 5\,s date realiste pentru nodul interior
(\texttt{esp32\_node\_a}): temperaturi de bază de 21\,°C în living și
19\,°C în dormitor, umiditate, gaz și mișcare aleatorie, inclusiv
câmpurile \texttt{temperature2}, \texttt{humidity2}, \texttt{light1},
\texttt{light2}, asigurând paritatea cu schema reală a payload-ului.
Când \texttt{SIMULATE=false}, \texttt{MqttService} preia controlul,
iar \texttt{SimulatorService} nu pornește timer-ul.

% ──────────────────────────────────────────────────────────────────
\section{Motorul de detecție a anomaliilor cu Isolation Forest}
\label{sec:motor-ml}

Detecția anomaliilor comportamentale constituie contribuția originală a
acestei lucrări.
Motorul \gls{ml} rulează în întregime local pe Raspberry Pi, fără
dependență de cloud, în acord cu principiul \emph{local-first}
descris în \chapref{ch:metodologie}~\cite{liu2008isolation,pedregosa2011scikit}.

\subsection{Ingineria trăsăturilor (\texttt{features.py})}
\label{sec:features}

Funcția \texttt{build\_features(df)} transformă un DataFrame Pandas cu
citiri brute în matricea de trăsături necesară modelului~\cite{mckinney2010pandas}.
Cele 12 trăsături, în ordinea fixă din \texttt{FEATURE\_NAMES}, sunt
descrise în \tabref{tab:features-cod}.

\begin{table}[H]
  \centering
  \caption{Cele 12 trăsături utilizate de motorul Isolation Forest.}
  \label{tab:features-cod}
  \begin{tabular}{@{} r l p{6.5cm} @{}}
    \toprule
    \textbf{Idx} & \textbf{Trăsătură} & \textbf{Sursă / transformare} \\
    \midrule
    0  & \texttt{temperature}      & DHT11 \#1, °C \\
    1  & \texttt{humidity}         & DHT11 \#1, \%RH \\
    2  & \texttt{gas\_normalized}  & MQ-2 ADC $/ 4095$ \\
    3  & \texttt{lux\_normalized}  & $\log_{1p}(\text{light1}) / 7{,}0$ (LDR \#1, scara ADC) \\
    4  & \texttt{motion}           & \gls{pir} 0/1 (de pe ESP32-CAM, GPIO13) \\
    5  & \texttt{hour\_of\_day}    & Ora din zi, 0–23 \\
    6  & \texttt{day\_of\_week}    & 0 (luni) – 6 (duminică) \\
    7  & \texttt{is\_weekend}      & 1 dacă zi = sâmbătă sau duminică \\
    8  & \texttt{is\_night}        & 1 dacă oră $\geq 23$ sau oră $< 6$ \\
    9  & \texttt{temp\_delta\_1h}  & Diferența de temperatură față de 2 citiri anterioare \\
    10 & \texttt{motion\_count\_10}& Suma evenimentelor PIR în fereastra de 10 citiri \\
    11 & \texttt{temp\_avg}        & Media DHT11 \#1 și \#2 (fallback = DHT11 \#1) \\
    \bottomrule
  \end{tabular}
\end{table}

Trăsăturile temporale (ora, ziua, weekend, noapte) permit modelului să
distingă comportamente normale specifice perioadei (ex. mișcare
nocturnă sau temperaturi ridicate în weekend) de adevăratele anomalii.
Normalizarea \texttt{StandardScaler} (medie zero, deviație standard unu)
este aplicată în \texttt{train.py} și reutilizată în \texttt{infer.py}.

\listingref{lst:features} prezintă calculul trăsăturilor derivate.

\begin{lstlisting}[language=Python,
  caption={Calcul trăsături derivate în \texttt{features.py} —
           normalizare gaz, lumină și media DHT11.},
  label=lst:features]
df['gas_normalized'] = df['gas_level'].fillna(0) / 4095.0

# Luminozitate: LDR ADC 0-1023; log1p/7 → interval ~[0, 1]
if 'light1' in df.columns:
    raw_light = df['light1'].fillna(
        df.get('light_lux', pd.Series(0, index=df.index)).fillna(0))
else:
    raw_light = df['light_lux'].fillna(0)
df['lux_normalized'] = np.log1p(raw_light) / 7.0

# Media DHT11 #1 si #2 (daca #2 lipseste, fallback la #1)
if 'temperature2' in df.columns:
    temp2 = df['temperature2'].fillna(df['temperature'])
else:
    temp2 = df['temperature']
df['temp_avg'] = (df['temperature'] + temp2) / 2.0
\end{lstlisting}

\subsection{Antrenarea modelului (\texttt{train.py})}
\label{sec:train}

Scriptul \texttt{train.py} extrage din PostgreSQL ultimele 30 de zile de
citiri fără alertă de gaz (citirile de gaz anormal ar perturba distribuția
normală), construiește matricea de trăsături, aplică
\texttt{StandardScaler} și antrenează un \texttt{IsolationForest} cu
parametrii: \texttt{contamination=0.05} (5\% din date asumate anomale),
\texttt{n\_estimators=100}, \texttt{random\_state=42}~\cite{liu2008isolation}.
Este impus un minim de 500 de rânduri; sub acest prag, scriptul se
oprește cu mesaj de informare, fără a suprascrie modelul existent.

Modelul antrenat și scalerul sunt persistate cu \texttt{joblib}
(\texttt{isolation\_forest.pkl}, \texttt{scaler.pkl}), iar metadatele
versiunii (număr de rânduri, contaminare, lista de trăsături) sunt
înregistrate în tabela \texttt{ml\_models} din PostgreSQL.
Reantrenarea este declanșată zilnic la 03:00 de serviciul systemd
\texttt{smarthome-ml.service} prin \texttt{scheduler.py}.

\subsection{Inferența în timp real (\texttt{infer.py})}
\label{sec:infer}

Clasa \texttt{AnomalyDetector} expune metoda \texttt{score(reading: dict)}
care primește o citire nouă (dicționar Python), construiește un DataFrame
de un singur rând, aplică \texttt{build\_features()} și
\texttt{scaler.transform()}, apoi calculează
\texttt{model.\allowbreak score\_samples(\allowbreak X\_scaled)[0]}.
Un scor sub pragul \texttt{ANOMALY\_THRESHOLD = -0.1} declanșează
crearea unei alerte de tip \texttt{ML\_ANOMALY} în baza de date,
urmată de notificarea aplicației mobile prin WebSocket~\cite{chandola2009anomaly}.

\textbf{Strategia cold-start}: în primele zile de funcționare, când
baza de date nu are suficiente date pentru antrenare ($< 500$ înregistrări),
motorul revine la reguli cu prag hardcodat (gaz $> 300$,
temperatură $> 35$\,°C sau $< 5$\,°C, mișcare în intervalul 23:00–06:00),
asigurând un nivel minim de protecție până când modelul statistic poate
fi inițializat.
Funcția \texttt{is\_ready()} permite codului apelant să distingă cele
două regimuri.

% ──────────────────────────────────────────────────────────────────
\section{Fluxuri principale — diagrame de secvență}
\label{sec:secvente}

Această secțiune descrie cele trei fluxuri principale ale sistemului
prin diagrame de secvență UML.

\subsection{Fluxul de monitorizare}
\label{sec:secv-monitor}

\figref{fig:secv-monitor} ilustrează ciclul de colectare și afișare a
datelor de senzori.
La fiecare 5\,s, ESP32 citește toți senzorii și publică un mesaj JSON pe
topicul \texttt{smarthome\slb{} esp32\_node\_a\slb{}sensors}.
\texttt{MqttService} (NestJS) recepționează mesajul, apelează
\texttt{SensorsService.saveReading()} care persistă datele în
TimescaleDB, emit evenimentul intern \texttt{sensor.reading}, iar
\texttt{SensorsGateway} transmite \texttt{sensor\_update} prin
Socket.IO tuturor clienților conectați la namespace-ul \texttt{/live}.
Aplicația mobilă actualizează dashboard-ul în timp real prin
hook-ul React Query care ascultă acest eveniment.

\begin{figure}[H]
  \centering
  \includegraphics[width=\singlefigure]{\pics/fig-secventa-monitorizare}
  \caption{Diagrama de secvență pentru fluxul de monitorizare:
           ESP32 → Mosquitto → NestJS → TimescaleDB → App.}
  \label{fig:secv-monitor}
\end{figure}

\subsection{Fluxul de comandă a actuatorilor}
\label{sec:secv-comanda}

\figref{fig:secv-comanda} descrie traseul unei comenzi de la utilizator
la actuatorul fizic.
Utilizatorul apasă butonul „LED 1 ON" în ecranul Control al aplicației,
care trimite \texttt{POST /api/commands/esp32\_node\_a} cu corpul
\texttt{\{"action":"led\_on","led":0\}}.
\texttt{CommandsController} (NestJS) procesează cererea, emite evenimentul
\texttt{command.sent}, iar \texttt{MqttService} publică comanda pe topicul
\texttt{smarthome\slb{} esp32\_node\_a\slb{}commands} (QoS 1).
ESP32 recepționează mesajul în callback-ul \texttt{onMessage()},
apelează \texttt{setLed(0, true)} și aplică nivel HIGH pe GPIO25.
Întreg fluxul se desfășoară în rețeaua locală, fără ieșire pe internet.

\begin{figure}[H]
  \centering
  \includegraphics[width=\singlefigure]{\pics/fig-secventa-comanda}
  \caption{Diagrama de secvență pentru fluxul de comandă a actuatorilor:
           App → NestJS → Mosquitto → ESP32 → actuator fizic.}
  \label{fig:secv-comanda}
\end{figure}

\subsection{Fluxul de detecție a anomaliilor}
\label{sec:secv-alerta}

\figref{fig:secv-alerta} prezintă traseul complet al unei anomalii
detectate de motorul \gls{if}.
O citire nouă sosită pe topicul \texttt{smarthome/+/sensors} este salvată
în baza de date, după care \texttt{SensorsService} apelează modulul ML
(\texttt{ml.service.ts}) care invocă sincron \texttt{AnomalyDetector.score()}.
Dacă scorul returnat este sub pragul $-0{,}1$, \texttt{AlertsService}
creează o înregistrare de tip \texttt{ML\_ANOMALY} în tabela
\texttt{alerts}, emit evenimentul \texttt{alert.new}, iar
\texttt{SensorsGateway} îl transmite aplicației prin Socket.IO.
Utilizatorul primește o notificare vizuală pe ecranul Alerte.

\begin{figure}[H]
  \centering
  \includegraphics[width=\singlefigure]{\pics/fig-secventa-alerta}
  \caption{Diagrama de secvență pentru detecția anomaliei comportamentale
           cu \gls{if}: NestJS apelează \texttt{infer.py} și generează
           alertă \texttt{ML\_ANOMALY}.}
  \label{fig:secv-alerta}
\end{figure}

% ──────────────────────────────────────────────────────────────────
\section{Aplicația mobilă}
\label{sec:mobile}

Interfața utilizatorului este implementată ca aplicație React Native 0.81
compilată cu Expo SDK 54, targetând iOS și Android~\cite{reactnative2024,expo2024}.
Starea globală este gestionată cu \textbf{Zustand}, apelurile \gls{rest}
cu \textbf{React Query}, iar datele în timp real prin \textbf{Socket.IO}
conectat la namespace-ul \texttt{/live} al \gls{api}-ului.
Autentificarea utilizează token \gls{jwt} stocat local.

\subsection{Discovery-ul gateway-ului}
\label{sec:discovery}

La prima pornire, aplicația încearcă să contacteze gateway-ul în
ordinea: variabila de mediu configurată la build-time,
\texttt{http://smarthome.local:3000} (mDNS) sau adresa fixă a
hotspot-ului de configurare (\texttt{192.168.4.1}).
Această strategie elimină nevoia configurării manuale a IP-ului în
scenariile tipice de rețea locală, susținând obiectivul de funcționare
„plug and play". Implementarea, prezentată în \listingref{lst:discovery},
sondează \emph{în paralel} toate adresele candidate și o reține pe prima
care răspunde la endpointul public \texttt{/api/setup/status}; la o eroare
de rețea adresa este invalidată, astfel încât aplicația redescoperă
gateway-ul atunci când telefonul trece, de exemplu, de pe hotspot pe
rețeaua Wi-Fi a casei.

\begin{lstlisting}[language=TypeScript,
  caption={Descoperirea automată a gateway-ului în rețea (\texttt{services/discovery.ts}).},
  label=lst:discovery]
const CANDIDATES = [
  API_HOST,                       // variabila de mediu / fallback
  'http://smarthome.local:3000',  // mDNS - orice retea
  HOTSPOT_HOST,                   // hotspot RPi - 192.168.4.1
];

// sondeaza in paralel; prima adresa care raspunde "castiga"
export async function resolveHost(force = false): Promise<string> {
  if (resolvedHost && !force) return resolvedHost;
  return Promise.any(CANDIDATES.map(probe))
    .then((host) => { resolvedHost = host; return host; })
    .catch(() => API_HOST);
}

async function probe(host: string): Promise<string> {
  // /api/setup/status e public si raspunde doar de pe RPi
  await axios.get(`${host}/api/setup/status`, { timeout: 4000 });
  return host;
}
\end{lstlisting}

\subsection{Recepția datelor în timp real (Socket.IO)}
\label{sec:mobile-ws}

Actualizarea live a interfeței nu se bazează pe interogare repetată
(\emph{polling}), ci pe o conexiune \textbf{Socket.IO} către namespace-ul
\texttt{/live} al \gls{api}-ului, autentificată cu același token \gls{jwt}.
La fiecare eveniment primit de la server, handler-ul actualizează direct
store-ul Zustand, iar componentele abonate se redesenează automat
(\listingref{lst:ws-client}). Reconectarea este automată, cu un număr
nelimitat de încercări — esențial pentru un dispozitiv mobil care își
schimbă frecvent starea rețelei.

\begin{lstlisting}[language=TypeScript,
  caption={Clientul Socket.IO: evenimente live $\rightarrow$ store global (\texttt{services/socket.ts}).},
  label=lst:ws-client]
socket = io(`${host}/live`, {
  auth: { token },
  transports: ['websocket'],
  reconnectionDelay: 3000,
  reconnectionAttempts: Infinity,
});

// citire noua de senzor -> actualizeaza store-ul -> UI live
socket.on('sensor_update', (reading: SensorReading) => {
  useAppStore.getState().setLatestReading(reading);
});
socket.on('alert', (alert: AlertItem) => {
  useAppStore.getState().addAlert(alert);          // badge + ecran Alerte
});
socket.on('node_status', (s: { nodeId: string; online: boolean }) => {
  useAppStore.getState().setNodeStatus(s.nodeId, s.online);
});
\end{lstlisting}

\subsection{Trimiterea comenzilor către actuatori}
\label{sec:mobile-cmd}

Comenzile către actuatori sunt expuse printr-o singură funcție generică
\texttt{api.\allowbreak commands.\allowbreak send(nodeId,\allowbreak action,\allowbreak extras)}, care construiește
corpul \gls{json} al cererii \gls{rest}; \texttt{extras} permite atașarea
parametrilor specifici (de exemplu indexul \gls{led}-ului sau unghiul
servomotorului). În ecranul de control, fiecare acțiune este înfășurată
într-o mutație React Query care aplică o \emph{actualizare optimistă} a
stării vizuale, astfel încât comutatorul reacționează instantaneu
(\listingref{lst:mobile-cmd}).

\begin{lstlisting}[language=TypeScript,
  caption={Trimiterea unei comenzi de actuator cu actualizare optimistă (\texttt{services/api.ts} + \texttt{ControlScreen.tsx}).},
  label=lst:mobile-cmd]
// services/api.ts - o singura functie pentru toate comenzile
send: (nodeId, action, extras) =>
  client.post(`/commands/${nodeId}`, { action, ...extras })
        .then((r) => r.data),

// ControlScreen.tsx - buton LED cu actualizare optimista
const mutation = useMutation({
  mutationFn: ({ action, led }) => api.commands.send(nodeId, action, { led }),
  onSuccess: (_, vars) =>
    setLeds((prev) => ({ ...prev, [vars.led]: vars.action === 'led_on' })),
});
// la apasarea butonului:
mutation.mutate({ action: on ? 'led_off' : 'led_on', led: led.index });
\end{lstlisting}

\subsection{Ecranele aplicației}
\label{sec:ecrane}

Aplicația este organizată în patru tab-uri principale accesibile după
autentificare, plus un ecran de setup:

\textbf{Dashboard} afișează un „Hero card" cu temperatura medie,
umiditatea, iluminarea și nivelul de gaz agregate de pe toate nodurile,
urmat de tile-uri individuale pentru fiecare nod.
Ecranul prezintă distinct datele celor doi senzori \gls{dht}11
(\texttt{temperature}/\texttt{humidity} pentru living și
\texttt{temperature2}/\texttt{humidity2} pentru dormitor)
și valorile brute ADC ale celor două \gls{ldr}-uri (\texttt{light1},
\texttt{light2}).
Sparkline-uri animate vizualizează tendința recentă a fiecărei valori,
iar un indicator „scor de siguranță" sintetizează starea generală.

\textbf{Control} permite acționarea tuturor actuatorilor:
\begin{itemize}
  \item Trei butoane individuale pentru \gls{led}-urile din living
        (GPIO25), dormitor (GPIO26) și hol (GPIO27);
  \item Comutator pentru ventilator (GPIO33 prin tranzistor);
  \item Trei preset-uri pentru servo (draperii): Deschis (0°), 50\%
        (90°), Închis (180°);
  \item Buton de test buzzer și buton „Oprire totală" (trimite comanda
        \texttt{all\_off});
  \item Scene predefinite: Acasă, Plec, Noapte, Cinema;
  \item Vizualizare cameră: snapshot (polling la 5\,s via
        \texttt{GET /capture}) și stream live MJPEG redat în
        \texttt{WebView} via \texttt{GET /stream}.
\end{itemize}

\textbf{Istoric} prezintă grafice de tendință pentru parametrii
selectați (temperatură, umiditate, gaz, lumină) pe intervale de
timp configurabile, folosind date din \texttt{GET /api/sensors/history}.

\textbf{Alerte} listează cronologic toate alertele generate de sistem
(\texttt{GAS\_DETECTED}, \texttt{MOTION\_NIGHT}, \texttt{ALERT\_HEAT},
\texttt{ALERT\_COLD}, \texttt{ML\_ANOMALY}), cu indicatori de severitate
și posibilitatea de a marca o alertă ca recunoscută
(\texttt{POST /api/alerts/\{id\}/acknowledge}).

\figref{fig:mobile} prezintă capturi reale din aplicație, ilustrând
limbajul vizual \emph{glassmorphism} adoptat: suprafețe translucide cu
estompare (\textit{blur}), accente vii pe gradient (indigo–cyan),
tipografie Space~Grotesk pentru valori numerice și Inter pentru text,
respectând contrastul \gls{ui} conform recomandărilor de accesibilitate.

\begin{figure}[H]
  \centering
  \subfloat[Dashboard]{\includegraphics[width=0.235\textwidth]{\pics/fig-mobile-dashboard}}\hfill
  \subfloat[Alerte]{\includegraphics[width=0.235\textwidth]{\pics/fig-mobile-alerte}}\hfill
  \subfloat[Control (sus)]{\includegraphics[width=0.235\textwidth]{\pics/fig-mobile-control1}}\hfill
  \subfloat[Control (jos)]{\includegraphics[width=0.235\textwidth]{\pics/fig-mobile-control2}}
  \caption{Capturi reale din aplicația mobilă: (a)~Dashboard — scor de
           siguranță, medii ambientale și tile-uri pe camere; (b)~Alerte —
           cronologie cu severitate și stare de confirmare; (c)–(d)~Control —
           cameră live, scene, draperii (servo), \gls{led}-uri și ventilator,
           cu indicator „nod offline" când nodul nu este conectat.}
  \label{fig:mobile}
\end{figure}

% ──────────────────────────────────────────────────────────────────
\section{Sumar}
\label{sec:implementare-sumar}

Capitolul a prezentat implementarea concretă a fiecărui strat al
sistemului smart home local-first.
Firmware-ul nodului interior (ESP32 WROOM-32) gestionează cei unsprezece
pini activi, publică un payload JSON standardizat la 5\,s și asigură
reconectarea automată la broker cu mecanisme \emph{Last Will} și
provisioning \gls{wifi} prin NVS.
ESP32-CAM adaugă capacitate video (stream MJPEG, snapshot JPEG) și
detecție de mișcare prin \gls{pir} GPIO13.
Gateway-ul Raspberry Pi orchestrează Mosquitto, PostgreSQL/TimescaleDB
și serviciile systemd.
\gls{api}-ul NestJS leagă toate componentele prin module clare (MQTT,
WebSocket, \gls{rest}, auth \gls{jwt}).
Contribuția originală — motorul Isolation Forest cu 12 trăsături —
rulează local și detectează anomalii comportamentale în timp real,
cu strategie cold-start pentru fazele inițiale de colectare a datelor.
Aplicația mobilă React Native/Expo oferă vizualizare live, control
complet al actuatorilor și acces la camera exterioară, toate fără
dependență de cloud.
